\documentclass[11pt,a4paper]{article}

\usepackage[UTF8]{ctex}
\usepackage{amsmath,amssymb}
\usepackage{geometry}
\usepackage{booktabs}
\usepackage{xcolor}
\geometry{margin=2.5cm}

\title{小车横向与纵向控制器推导}
\author{}
\date{\today}

\begin{document}
\maketitle

\section{车辆模型与运动学}

阿克曼转向小车，轴距 $L = 0.15\,\text{m}$，后轮距 $w = 0.15\,\text{m}$。

\subsection{自行车运动学}

\begin{align}
\dot{x}   &= v \cos\theta \\
\dot{y}   &= v \sin\theta \\
\dot\theta &= \frac{v}{L}\tan\delta \equiv \omega_{\text{bicycle}} \label{eq:bicycle}
\end{align}

\subsection{误差动力学 (虚拟车参考系)}

定义误差：
\begin{align}
e_x &= (x_{\text{vst}} - x_{\text{car}})\cos\theta_{\text{vst}} + (y_{\text{vst}} - y_{\text{car}})\sin\theta_{\text{vst}} \label{eq:ex}\\
e_y &= (x_{\text{car}} - x_{\text{vst}})\sin\theta_{\text{vst}} - (y_{\text{car}} - y_{\text{vst}})\cos\theta_{\text{vst}} \label{eq:ey}\\
e_\theta &= \theta_{\text{vst}} - \theta_{\text{car}} \\
e_\delta &= \delta_{\text{vst}} - \delta_{\text{car}}
\end{align}

\begin{center}
\begin{tabular}{c|l}
\toprule
符号 & 名称 \\
\midrule
$e_x$ & 纵向误差 \\
$e_y$ & 横向误差 \\
$e_\theta$ & 航向误差 \\
$\dot{e}_\theta$ & 横摆角速度误差 \\
$e_\delta$ & 舵角误差 \\
$e_y^{\text{int}}$ & 横向误差积分 \\
\bottomrule
\end{tabular}
\end{center}

小角度线性化下 ($\sin e_\theta \approx e_\theta$)：
\begin{align}
\dot{e}_y &= v \cdot e_\theta \label{eq:ey_dot}\\
\dot{e}_\theta &= \omega_{\text{vst}} - \omega_{\text{car}} \equiv \dot{e}_\theta
\end{align}

\section{横向控制器: LQR}

\subsection{被控对象: 一阶惯性 augmentation}

自行车模型隐含假设：轮胎力瞬时响应。现实中轮胎滑移、舵机柔度等导致横摆角速度有一阶延迟。建模为：

\begin{equation}
\ddot{e}_\theta = -w_o \cdot \dot{e}_\theta + w_o \cdot \frac{v}{L} \cdot e_\delta
\label{eq:ethd_dyn}
\end{equation}

其中 $w_o$ 为虚拟阻尼带宽 (当前值 $w_o = 15\,\text{rad/s}$)。

传递函数形式：
\begin{equation}
\dot{e}_\theta = \frac{w_o}{s + w_o} \cdot \frac{v}{L} \cdot e_\delta
\end{equation}

稳态 ($s \to 0$) 退化为纯自行车模型 $\dot{e}_\theta = (v/L)e_\delta$。

\subsection{增广状态空间 (5-term LQI)}

状态向量 $\mathbf{x} = [e_y^{\text{int}},\; e_y,\; e_\theta,\; \dot{e}_\theta,\; e_\delta]^\top$，控制输入 $u = \omega$ (转向角速度)。5 个状态对应 5 个 Riccati 增益，无需分解。

系统矩阵：
\begin{equation}
\mathbf{A} =
\begin{bmatrix}
0 & 1 & 0 & 0      & 0            \\
0 & 0 & v & 0      & 0            \\
0 & 0 & 0 & 1      & 0            \\
0 & 0 & 0 & -w_o   & w_o \cdot v/L \\
0 & 0 & 0 & 0      & 0
\end{bmatrix},
\quad
\mathbf{B} =
\begin{bmatrix}
0 \\ 0 \\ 0 \\ 0 \\ 1
\end{bmatrix}
\label{eq:AB}
\end{equation}

各行含义：
\begin{align}
\dot{e}_y^{\text{int}} &= e_y                            &&\text{积分器} \\
\dot{e}_y              &= v \cdot e_\theta                &&\text{运动学耦合} \\
\dot{e}_\theta         &= \dot{e}_\theta                  &&\text{定义} \\
\ddot{e}_\theta        &= -w_o \dot{e}_\theta + \frac{w_o v}{L} e_\delta &&\text{一阶惯性} \\
\dot{e}_\delta         &= \omega                          &&\text{舵机积分}
\end{align}

\subsection{代价函数}

\begin{equation}
J = \int_0^\infty \big(
q_{\text{int}} e_y^{\text{int}\,2}
+ q_{ey} e_y^2
+ q_{\text{eth}} e_\theta^2
+ q_{\text{ethd}} \dot{e}_\theta^2
+ r\,\omega^2
\big) dt
\label{eq:cost}
\end{equation}

当前参数: $q_{\text{int}} = 0.05$, $q_{ey} = 25000$, $q_{\text{eth}} = 1000$, $q_{\text{ethd}} = 35$, $r = 1$。

\subsection{Riccati 方程求解}

连续代数 Riccati 方程：
\begin{equation}
\mathbf{A}^\top \mathbf{P} + \mathbf{P}\mathbf{A} - \mathbf{P}\mathbf{B}r^{-1}\mathbf{B}^\top\mathbf{P} + \mathbf{Q} = 0
\end{equation}

其中 $\mathbf{Q} = \operatorname{diag}(q_{\text{int}},\, q_{ey},\, q_{\text{eth}},\, q_{\text{ethd}},\, 0)$。

最优反馈增益：
\begin{equation}
\mathbf{K} = r^{-1}\mathbf{B}^\top\mathbf{P} = [K_i,\; K_{ey},\; K_{\text{eth}},\; K_{\text{ethd}},\; K_{ed}]
\end{equation}

\subsection{控制律 (5-term)}

\begin{equation}
\boxed{
\omega^{\text{fb}} = K_i e_y^{\text{int}} + K_{ey} e_y + K_{\text{eth}} e_\theta + K_{\text{ethd}} \dot{e}_\theta + K_{ed} e_\delta
}
\label{eq:omega_fb}
\end{equation}

\begin{equation}
\omega_{\text{cmd}} = \omega^{\text{ff}} + \omega^{\text{fb}}
\label{eq:omega_cmd}
\end{equation}

其中 $\omega^{\text{ff}} = \dot{\delta}_{\text{vst}}$ 为模型前馈。

舵机指令：
\begin{equation}
\delta_{\text{cmd}} = \delta_{\text{car}} + \omega_{\text{cmd}} \cdot \Delta t
\end{equation}

约束在 $|\delta| \leq \delta_{\max} = 0.4636\,\text{rad}\,(26.5^\circ)$, $|\omega| \leq \omega_{\max} = 20\,\text{rad/s}$。

\subsection{反馈状态来源}

\begin{center}
\begin{tabular}{c|c|c|c}
\toprule
状态 & 表达式 & 真实值来源 & 模型/参考来源 \\
\midrule
$e_y^{\text{int}}$ & $\int e_y \,dt$ & --- (纯计算) & --- \\
$e_y$ & 式 (\ref{eq:ey}) & 航位推算 (IMU+编码器) & vst 运动学积分 \\
$e_\theta$ & $\theta_{\text{vst}} - \theta_{\text{car}}$ & IMU 四元数 (纯, 不融合陀螺) & vst 运动学积分 \\
$\dot{e}_\theta$ & $\omega_{\text{bicycle}}(vst) - \omega_z^{\text{gyro}}$ & IMU 陀螺仪 z 轴 & 自行车模型 \\
$e_\delta$ & $\delta_{\text{vst}} - \delta_{\text{car}}$ & 上周期舵机指令 & vst 运动学积分 \\
\bottomrule
\end{tabular}
\end{center}

\textbf{关键设计}: $\dot{e}_\theta$ 使用 IMU 陀螺仪直接测量而非模型预测，$e_\theta$ 使用 IMU 四元数直接解算。LQR 设计时用 $w_o$ 一阶惯性对 $\dot{e}_\theta$ 建模，但代码实现用更真实的传感器替代，反馈更及时。

\subsection{积分反饱和}

\begin{equation}
e_y^{\text{int}} \leftarrow \operatorname{clamp}\!\left(e_y^{\text{int}} + e_y \cdot \Delta t,\; -e_{y,\max}^{\text{int}},\; +e_{y,\max}^{\text{int}}\right), \quad e_{y,\max}^{\text{int}} = 0.5 \; [\text{m}\!\cdot\!\text{s}]
\end{equation}

\subsection{虚拟车状态同步 (安全网)}

当位置误差 $>5\,\text{cm}$ 时，vst 全状态同步至真实车：
\begin{equation}
(x_{\text{vst}},\, y_{\text{vst}},\, \theta_{\text{vst}},\, v_{\text{vst}},\, \delta_{\text{vst}}) \leftarrow (x_{\text{car}},\, y_{\text{car}},\, \theta_{\text{car}},\, v_{\text{car}},\, \delta_{\text{car}})
\end{equation}

防止模型失配导致 vst 漂移不可恢复。

\section{纵向控制器: 降阶 LADRC}

\subsection{纵向动力学模型}

采用一阶阻力模型：
\begin{equation}
\dot{v} = -\alpha \cdot v + b_0 \cdot u + f
\label{eq:long_model}
\end{equation}

其中 $\alpha = 1.5$ (模型阻尼)，$b_0 = 30$ (控制增益)，$u \in [-1, 1]$ (归一化油门)，$f$ 为总扰动 (坡度、风阻、摩擦等)。

稳态特性: $v_{ss} = (b_0/\alpha) \cdot u \approx 20 \cdot u$。

\subsection{降阶 LESO (仅观测扰动)}

全阶 LESO 对无传感器延迟系统是冗余的——速度 $v$ 已由编码器测量。降阶观测器仅估计扰动 $f$。

定义辅助变量：
\begin{equation}
z = \hat{f} - w_o \cdot v
\label{eq:z_def}
\end{equation}

扰动常值假设 ($\dot{f} \equiv 0$) 下的动态：
\begin{align}
\dot{z} &= \dot{\hat{f}} - w_o \cdot \dot{v} \\
        &= 0 - w_o \cdot (-\alpha v + b_0 u + \hat{f}) \\
        &= -w_o \cdot \big(z + (w_o - \alpha)v + b_0 u\big)
\label{eq:zdot}
\end{align}

离散化 (前向欧拉)：
\begin{align}
z_{k+1} &= z_k + \dot{z} \cdot \Delta t \\
\hat{f}_{k} &= z_k + w_o \cdot v_{\text{meas}}
\end{align}

参数: $w_o = 15\,\text{rad/s}$ (观测带宽)，$\Delta t = 5\,\text{ms}$ (200Hz)。

\subsection{PD 控制律}

\begin{equation}
u_0 = k_p \cdot e_x + k_d \cdot (v_{\text{ref}} - v_{\text{meas}}) + a_{\text{ref}} + \alpha \cdot v_{\text{meas}} - \hat{f}
\label{eq:u0}
\end{equation}

\begin{equation}
u = \operatorname{clamp}\!\left(\frac{u_0}{b_0},\; -1,\; +1\right)
\label{eq:u}
\end{equation}

其中 $k_p = 16$ (位置增益)，$k_d = 8$ (速度阻尼，$\zeta = 1.0$)，$a_{\text{ref}}$ 为 NN 前馈加速度。

控制律结构解读：
\begin{itemize}
\item $k_p \cdot e_x$: 纵向位置误差的 PD 修正 ($e_x$ 来自式 (\ref{eq:ex}))
\item $k_d \cdot (v_{\text{ref}} - v_{\text{meas}})$: 速度阻尼
\item $a_{\text{ref}}$: NN 前馈加速度
\item $\alpha \cdot v_{\text{meas}}$: 模型阻尼补偿 (抵消 $-\alpha v$ 项)
\item $-\hat{f}$: 扰动补偿 (抵消 $f$)
\end{itemize}

理想补偿后，闭环退化为：
\begin{equation}
\dot{v} \approx k_p \cdot e_x + k_d \cdot (v_{\text{ref}} - v_{\text{meas}}) + a_{\text{ref}}
\end{equation}

即位置-速度双环 PD + 前馈。

\subsection{阿克曼差速}

左右轮速分配：
\begin{align}
u_L &= u \cdot \left(1 - \frac{w}{2} \cdot \frac{\tan\delta}{L}\right) \\
u_R &= u \cdot \left(1 + \frac{w}{2} \cdot \frac{\tan\delta}{L}\right)
\end{align}

其中 $w = 0.15\,\text{m}$ 为后轮距。$|\delta| < 10^{-4}$ 时退化为 $u_L = u_R = u$。

\section{完整控制器框图}

\[
\begin{array}{c}
\text{NN 规划器 (20Hz)} \\
\downarrow \quad (a_{\text{ref}},\; \omega_{\text{ff}}) \\
\begin{array}{|c|}
\hline
\text{虚拟车运动学 (200Hz)} \\
\dot{x} = v\cos\theta,\; \dot{y} = v\sin\theta,\; \dot\theta = \omega_{\text{bicycle}},\; \dot{v} = a_{\text{ref}} \\
\hline
\end{array} \\
\downarrow \quad (x_{\text{vst}},\; y_{\text{vst}},\; \theta_{\text{vst}},\; v_{\text{vst}},\; \delta_{\text{vst}}) \\
\begin{array}{|c|c|}
\hline
\textbf{横向 LQR (200Hz)} & \textbf{纵向 LADRC (200Hz)} \\
\mathbf{K}_{5\times1} \cdot \mathbf{x} & \text{降阶 LESO + PD} \\
\omega_{\text{cmd}} \to \delta_{\text{cmd}} & u \to (u_L, u_R) \\
\hline
\end{array} \\
\downarrow \\
\begin{array}{|c|}
\hline
\text{执行器} \\
\text{舵机}\; \delta,\; \text{电机} \\
\hline
\end{array} \\
\downarrow \\
\begin{array}{|c|}
\hline
\text{真实车辆 + 传感器} \\
\text{IMU (}\theta,\; \omega_z\text{), 编码器 (}v\text{), 舵机指令 (}\delta\text{)} \\
\hline
\end{array} \\
\downarrow \quad (x_{\text{car}},\; y_{\text{car}},\; \theta_{\text{car}},\; v_{\text{car}},\; \delta_{\text{car}}) \\
\begin{array}{|c|}
\hline
\text{误差投影 } e_x, e_y, e_\theta, \dot{e}_\theta, e_\delta \\
\hline
\end{array} \\
\uparrow \quad \text{反馈至 LQR/LADRC}
\end{array}
\]

\section{参数总表}

\begin{center}
\begin{tabular}{l|c|c}
\toprule
参数 & 符号 & 当前值 \\
\midrule
\multicolumn{3}{c}{\textbf{车体}} \\
轴距 & $L$ & 0.15 m \\
后轮距 & $w$ & 0.15 m \\
舵角限幅 & $\delta_{\max}$ & 0.4636 rad (26.5°) \\
舵速限幅 & $\omega_{\max}$ & 20 rad/s \\
\midrule
\multicolumn{3}{c}{\textbf{横向 LQR}} \\
ey 积分权重 & $q_{\text{int}}$ & 0.05 \\
ey 权重 & $q_{ey}$ & 25000 \\
$e_\theta$ 权重 & $q_{\text{eth}}$ & 1000 \\
$\dot{e}_\theta$ 权重 & $q_{\text{ethd}}$ & 35 \\
控制权重 & $r$ & 1 \\
虚拟阻尼 & $w_o$ & 15 rad/s \\
积分反饱和上限 & $e_{y,\max}^{\text{int}}$ & 0.5 m·s \\
vst 同步门限 & --- & 0.05 m \\
速度断点 & --- & 0.1--5.0 m/s (11 点) \\
\midrule
\multicolumn{3}{c}{\textbf{纵向 LADRC}} \\
控制增益 & $b_0$ & 30 \\
模型阻尼 & $\alpha$ & 1.5 \\
观测带宽 & $w_o$ & 15 rad/s \\
位置增益 & $k_p$ & 16 \\
速度增益 & $k_d$ & 8 \\
控制频率 & $1/\Delta t$ & 200 Hz \\
\bottomrule
\end{tabular}
\end{center}

\end{document}
